\documentclass{article}
\usepackage[utf8]{inputenc}
\usepackage{amsmath, amssymb, amsfonts}
\usepackage{algorithm}
\usepackage{algpseudocode}
\usepackage{listings}
\usepackage{xcolor}
\usepackage{geometry}
\geometry{a4paper, margin=1in}

% Python code styling for listings
\definecolor{codegreen}{rgb}{0,0.6,0}
\definecolor{codegray}{rgb}{0.5,0.5,0.5}
\definecolor{codepurple}{rgb}{0.58,0,0.82}
\definecolor{backcolour}{rgb}{0.95,0.95,0.92}

\lstdefinestyle{mystyle}{
    backgroundcolor=\color{backcolour},   
    commentstyle=\color{codegreen},
    keywordstyle=\color{magenta},
    numberstyle=\tiny\color{codegray},
    stringstyle=\color{codepurple},
    basicstyle=\ttfamily\footnotesize,
    breakatwhitespace=false,         
    breaklines=true,                 
    captionpos=b,                    
    keepspaces=true,                 
    numbers=left,                    
    numbersep=5pt,                  
    showspaces=false,                
    showstringspaces=false,
    showtabs=false,                  
    tabsize=2
}
\lstset{style=mystyle}

\title{Matrix-Free Krylov Solvers for Large Adjoint Systems in JAX}
\author{}
\date{}

\begin{document}

\maketitle

\section{Problem Formulation}

We consider a large-scale discrete linear system structured over time steps $t = 1, \dots, T$. The governing equation is given by:
\begin{equation}
    A_i \lambda_i + A_{i+1} \lambda_{i+1} = b_i, \quad \forall i \in \{1, \dots, T-1\}
    \label{eq:system}
\end{equation}
where:
\begin{itemize}
    \item $\lambda_i \in \mathbb{R}^N$ is the unknown state vector at step $i$.
    \item $b_i \in \mathbb{R}^N$ is the forcing term.
    \item $A_i \in \mathbb{R}^{N \times N}$ is the Jacobian matrix at step $i$, typically defined as $A_i = \frac{\partial f}{\partial x}(x_i)$ for some dynamics function $f$.
\end{itemize}

\subsection{Constraint: Memory Efficiency}
For large $N$ and $T$, explicitly storing the matrices $A_i$ is computationally prohibitive ($O(T \cdot N^2)$ memory). We require an implementation that solves Eq. \eqref{eq:system} using $O(T \cdot N)$ memory by leveraging \textbf{Matrix-Free Krylov methods}.

\section{JAX Implementation Strategy}

The solution utilizes \texttt{jax.scipy.sparse.linalg.gmres}. Instead of passing a materialized matrix, we pass a linear operator function $\mathcal{L}(\lambda)$ that computes the left-hand side of Eq. \eqref{eq:system} on the fly.

\subsection{Forward Mode: The Jacobian-Vector Product (JVP)}
If $A_i$ represents the standard Jacobian, we compute the product $A_i v$ using Forward-Mode Automatic Differentiation via \texttt{jax.jvp}.

Given a primal trajectory $X = \{x_1, \dots, x_T\}$, the operation is:
\begin{equation}
    A_i \lambda_i \equiv \texttt{jvp}(f, x_i, \lambda_i)
\end{equation}
This operation computes the product in a single pass without instantiating the Jacobian matrix.

\subsection{Reverse Mode: The Vector-Jacobian Product (VJP)}
In many adjoint formulations (e.g., backpropagation or sensitivity analysis), the system often involves the \textbf{transpose} $A_i^\top \lambda_i$. In this case, we use Reverse-Mode Automatic Differentiation via \texttt{jax.vjp}:
\begin{equation}
    A_i^\top \lambda_i \equiv \texttt{vjp}(f, x_i)(\lambda_i)
\end{equation}

\section{Algorithm and Code}

The following JAX implementation solves the system without pre-computing matrices. It vectorizes the operation over the time axis using \texttt{jax.vmap}.

\begin{lstlisting}[language=Python, caption=Matrix-Free GMRES Implementation]
import jax
import jax.numpy as jnp
from jax.scipy.sparse.linalg import gmres

def solve_adjoint_system(x_trajectory, b_rhs, f_dynamics):
    """
    Solves A_i * lam_i + A_{i+1} * lam_{i+1} = b_i
    where A_i = Jacobian(f_dynamics) at x_trajectory[i]
    """
    
    T, N = x_trajectory.shape

    # Define the Matrix-Free Linear Operator
    def matvec_operator(lam):
        # 1. Helper to compute A * v using JVP (Forward Mode)
        def apply_jacobian(x, v):
            # jvp returns (primals_out, tangents_out)
            # tangents_out is exactly J @ v
            _, Jv = jax.jvp(f_dynamics, (x,), (v,))
            return Jv

        # 2. Term 1: A_i * lam_i
        # Vectorize over time dimension (axis 0)
        term1 = jax.vmap(apply_jacobian)(x_trajectory, lam)

        # 3. Term 2: A_{i+1} * lam_{i+1}
        # Shift lambda and x forward by 1 step to align indices
        lam_shifted = jnp.roll(lam, shift=-1, axis=0)
        lam_shifted = lam_shifted.at[-1].set(jnp.zeros(N)) # Boundary condition
        
        x_shifted = jnp.roll(x_trajectory, shift=-1, axis=0)
        x_shifted = x_shifted.at[-1].set(jnp.zeros(N))

        term2 = jax.vmap(apply_jacobian)(x_shifted, lam_shifted)

        return term1 + term2

    # Solve using GMRES
    # JAX traces the operator and optimizes the JVP computations
    lam_sol, info = gmres(
        matvec_operator, 
        b_rhs, 
        tol=1e-5, 
        maxiter=100
    )
    
    return lam_sol, info

# Example Usage
if __name__ == "__main__":
    # Dummy dynamics: f(x) = sin(x)
    f = lambda x: jnp.sin(x)
    
    # Random data
    key = jax.random.PRNGKey(0)
    X = jax.random.normal(key, (100, 50)) # T=100, N=50
    B = jax.random.normal(key, (100, 50))
    
    sol, info = solve_adjoint_system(X, B, f)
    print(f"Convergence Info: {info}")
\end{lstlisting}

\section{Summary}
\begin{itemize}
    \item \textbf{Memory Complexity:} $O(TN)$ (vs $O(TN^2)$ for explicit matrices).
    \item \textbf{Compute Complexity:} Dependent on the cost of evaluating $f(x)$. The overhead of JVP is a small constant factor over the forward pass.
    \item \textbf{Preconditioning:} If convergence is slow, a block-diagonal preconditioner can be implemented by approximating $M \approx \text{diag}(A_1, \dots, A_T)$ and inverting blocks locally.
\end{itemize}

\end{document}